
% =============================================================================
% Section 01 — The Ontological Problem of Unity and Plurality
% Book I — Ontology of Viability
% Chapter 03 — Unity and Plurality
%
% Purpose:
%   To formulate the ontological problem of unity and plurality,
%   identify its principal conceptual dimensions, and establish
%   the distinctions required for its subsequent examination.
%
% Conceptual Policy:
%   1. Unity and plurality are treated as ontological problems,
%      not merely logical, linguistic, or epistemic problems.
%   2. Unity is not initially reduced to numerical singularity,
%      aggregation, structural coherence, or invariance.
%   3. Plurality is not initially reduced to plurality of parts,
%      manifestations, relations, or identities.
%   4. The relation between identity and unity remains open
%      for explicit examination in the following sections.
%   5. No assumption is made that every existent has exactly
%      one identity or that plurality necessarily entails
%      multiple identities.
% =============================================================================

\section{The Ontological Problem of Unity and Plurality}
\label{sec:ontological-problem-of-unity-and-plurality}

\subsection{The Question of Unity and Plurality}

The relation between unity and plurality constitutes a fundamental
problem in ontology. Every account of existence that encounters
determinate existents must also confront the question of how
something is one, how multiplicity is possible, and how unity
and plurality can belong to the same ontological field without
being confused or prematurely reduced to one another.

The problem is not exhausted by the observation that reality
contains many things. Nor is it resolved by asserting that each
thing is, in some respect, one. Such statements identify the
presence of unity and plurality but do not explain their
ontological status or their relation.

The central question is therefore not simply whether unity and
plurality exist, but what it means for an existent to possess
unity, what it means for multiplicity to obtain, and under what
conditions these determinations can belong to the same existent
or to related existents.

An existent may be considered one while also exhibiting internal
differentiation, multiple determinations, distinct components,
various relations, or several modes of manifestation. Conversely,
a plurality of existents may be considered together as a group,
collection, or organized whole. These situations are not
ontologically identical. They involve different senses of unity
and plurality, and their differences require explicit analysis.

The first task of this chapter is consequently to prevent the
problem from being reduced to a single relation between one
and many. Unity and plurality are not sufficiently understood
through a numerical opposition alone. Their ontological
significance depends on what is unified, what is multiple,
and the respect in which either determination is asserted.

\subsection{Unity as an Ontological Determination}

Unity cannot be treated merely as a numerical description
superimposed upon an existent from an external standpoint.
The question of unity concerns the ontological determination
of an existent as one, rather than only the manner in which
it is counted, named, classified, or considered.

To call something one may indicate different claims:

\begin{itemize}
    \item that an existent possesses a unified identity;
    \item that several components belong to one organized whole;
    \item that multiple determinations belong to one existent;
    \item that a plurality is collected under a common designation;
    \item that an existent is numerically distinct from another;
    \item that a configuration exhibits a certain degree of
    structural coherence.
\end{itemize}

These claims must not be assumed to express one and the same
ontological condition.

A collection can be one collection without necessarily
constituting one identity. An organized whole can possess
structural unity without that unity being identical to the
identity of its components or to the identity of the whole.
Likewise, the numerical distinction of an existent from another
does not, by itself, explain the ontological constitution of
its unity.

Unity must therefore be investigated according to the
ontological determination to which it refers. Its meaning
cannot be established solely by the grammatical or numerical
form of a statement.

The present section does not yet provide a complete definition
of unity. It establishes that unity is a problem requiring
differentiation among several possible ontological senses.

\subsection{Plurality as an Ontological Determination}

Plurality is equally resistant to a single undifferentiated
interpretation.

The presence of more than one component does not necessarily
establish the presence of more than one identity. The presence
of several relations does not necessarily establish several
existents. The multiplicity of manifestations does not,
without further argument, establish a multiplicity of the
identities manifested.

Plurality may concern different ontological determinations,
including:

\begin{itemize}
    \item plurality of existents;
    \item plurality of identities;
    \item plurality of components;
    \item plurality of determinations;
    \item plurality of relations;
    \item plurality of manifestations;
    \item plurality of configurations or compositions.
\end{itemize}

These forms of plurality cannot be treated as interchangeable.
Each concerns a different possible object of multiplicity
and requires its own ontological examination.

In particular, plurality of components and plurality of
identities must remain distinct questions. A composed existent
may contain multiple components without the number of its
components determining the number of identities involved.
Similarly, the multiplicity of manifestations associated with
an identity cannot be converted into a plurality of identities
without establishing an additional ontological criterion.

Plurality must therefore be specified according to what is
multiple. An assertion that something is multiple remains
ontologically incomplete until the relevant multiplicity
and its domain have been identified.

\subsection{The Apparent Opposition Between the One and the Many}

Unity and plurality are frequently presented as opposing
determinations. If something is one, it may appear that it
cannot also be many; if something is many, it may appear
that it cannot possess unity.

This opposition, however, is not yet an established ontological
contradiction. It may result from treating distinct senses
of unity and plurality as if they referred to the same
determination in the same respect.

For example, the unity of an existent and the plurality
of its components are not necessarily contradictory claims.
The first concerns the existent as a unified ontological
referent; the second concerns the multiplicity of elements
associated with its composition.

Likewise, the unity of an identity and the plurality of its
manifestations cannot be declared contradictory merely
because the terms one and many are applied within the same
ontological situation. The relevant question is whether
they refer to the same ontological subject and the same
determination, or to distinct but related determinations.

This distinction does not automatically resolve every
problem concerning unity and plurality. It only establishes
that an apparent opposition cannot be treated as a genuine
contradiction without first specifying the terms of the
opposition.

A valid ontological analysis must therefore distinguish
between:

\begin{enumerate}
    \item a genuine incompatibility between two determinations;
    \item the coexistence of distinct determinations within
    one existent;
    \item the relation between a unified existent and a plurality
    associated with it;
    \item the unity attributed to a plurality considered as
    a collection or organized whole.
\end{enumerate}

The distinction between these situations is a condition
of the inquiry, not its final conclusion.

\subsection{The Problem of the Subject of Unity}

Any claim concerning unity requires clarification of its
subject. The expression ``one'' does not identify, by itself,
what possesses unity or what kind of unity is being asserted.

The subject of unity may be an existent, an identity, a
composition, a structure, a relation, or an aggregate.
These subjects are not interchangeable.

The unity of an existent must not be assumed to be identical
to the unity of its composition. Nor can the unity of a
structure automatically be identified with the unity of
the existent whose determinations it organizes.

The same difficulty arises when plurality is considered.
A plurality may belong to an existent as its components,
to an identity through its manifestations, or to a domain
through the presence of several distinct existents.

Consequently, the ontological analysis of unity and plurality
requires two questions to be kept separate:

\begin{enumerate}
    \item What is the subject to which unity or plurality
    is attributed?
    \item What ontological determination is being described
    as unified or multiple?
\end{enumerate}

Failure to distinguish these questions leads to the
unwarranted transfer of properties from one ontological
subject to another.

The unity of an aggregate, for example, cannot establish
the unity of an identity merely because the aggregate
is considered one object. Likewise, the plurality of
components cannot establish the plurality of identities
merely because several components are present.

\subsection{Unity, Identity, and Numerical Distinction}

The problem of unity is closely related to identity,
but the two concepts must not be identified without
argument.

The preceding chapter established that invariance does
not constitute identity and that the preservation of
invariant determinations cannot independently establish
numerical identity. The present inquiry must preserve
the same conceptual discipline: unity cannot be assumed
to provide a complete definition or criterion of identity.

The unity of identity concerns the identity as one.
The numerical distinction between identities concerns
whether two identity-referents are numerically distinct.
The unity of a composition concerns the organization
or belonging together of its constituents.

These questions are related, but they are not equivalent.

In particular, the statement that an identity is one
does not yet explain whether an existent can possess
multiple identities, whether distinct identities can
belong to one composed existent, or whether one identity
can be associated with multiple manifestations.

Those questions require separate examination.

The present section therefore adopts no unrestricted
principle according to which every existent must possess
exactly one identity. It also does not assume that
multiple identities can be attributed to one existent
without qualification. Both positions require further
ontological argument.

\subsection{Unity and Singularity}

Unity must also be distinguished from singularity.

Singularity may refer to numerical uniqueness, the
non-duplication of an existent, or its distinction from
other existents. Unity, by contrast, concerns the
ontological condition under which something is one.

Although these concepts may coincide in particular
cases, their meanings are not identical.

An existent may be considered one without the concept
of unity thereby explaining its numerical uniqueness.
Conversely, a plurality may possess some form of unity
without becoming numerically singular in every relevant
respect.

The fact that a collection is treated as one collection,
for example, does not make its components numerically
singular. Similarly, the numerical distinctness of an
existent does not by itself explain the unity of its
identity or composition.

The distinction is essential because an account that
reduces unity to singularity risks treating every question
of unity as a question of counting. Such a reduction
would leave unexplained the ontological relations among
identity, composition, multiplicity, and the unity of
organized wholes.

The non-equivalence of unity and singularity will be
examined directly in Section~\ref{sec:non-equivalence-of-unity-and-singularity}.

\subsection{Unity and Composition}

Composition introduces a further dimension of the problem.

An existent may possess components, and those components
may stand in determinate relations to one another.
Their organization may constitute a composition or
an organized whole.

However, the unity of a composition cannot be assumed
to be identical to the unity of an identity.

The fact that several components belong to one composition
does not establish that they constitute one identity in
every ontological respect. Equally, the fact that a
composition undergoes changes in its components does
not, without further argument, determine whether its
identity remains, changes, or ceases.

The relation between unity and composition must therefore
be investigated without presupposing that the whole is
merely the sum of its parts or that the unity of the
whole is reducible to the relations among those parts.

This inquiry also requires distinguishing the unity
of an aggregate from the unity of an existent whose
composition is ontologically constitutive of it.

The distinction between these possibilities will be
developed in Sections~\ref{sec:unity-of-source-and-aggregate}
and~\ref{sec:plurality-and-composition}.

\subsection{The Problem of Plurality Within and Beyond an Existent}

Plurality can be considered internally or externally,
but this distinction must not be mistaken for a complete
classification of multiplicity.

Internal plurality concerns multiplicity attributed
to or present within an existent. It may involve
components, determinations, relations, or other
ontological differentiations.

External plurality concerns the presence of several
distinct existents within a wider domain or relation.

Neither category, by itself, determines the number
of identities involved.

Internal plurality does not necessarily entail
multiple identities, while external plurality
cannot be adequately described merely by counting
components or manifestations.

Moreover, the distinction between internal and external
plurality depends on the ontological boundaries and
determinations under examination. These boundaries
must not be assumed to be self-evident in every case.

The analysis must therefore remain attentive to the
difference between:

\begin{itemize}
    \item multiplicity within one existent;
    \item multiplicity among distinct existents;
    \item multiplicity among identities;
    \item multiplicity among manifestations;
    \item multiplicity generated or constituted through
    composition and relation.
\end{itemize}

The following sections will examine these distinctions
without assuming that one category can be reduced
to another.

\subsection{The Ontological Stakes of the Problem}

The problem of unity and plurality affects several
fundamental questions of ontology.

First, it concerns the status of identity. If unity
is reduced to aggregation, the identity of an existent
risks being treated as nothing more than the collection
of its constituents. If unity is reduced to structure,
identity risks being identified with a particular
organization of determinations. Neither reduction
can be accepted without an independent argument.

Second, the problem concerns the status of composition.
If plurality is treated as incompatible with unity,
the ontological possibility of composed existents
becomes difficult to explain. If unity is attributed
to every collection without qualification, however,
the distinction between an aggregate and an ontologically
unified existent disappears.

Third, the problem concerns the distinction among
different identities. If plurality of components,
relations, or manifestations is automatically interpreted
as plurality of identities, the number of identities
will be determined by categories that do not necessarily
correspond to identity itself.

Fourth, the problem concerns the limits of ontological
claims. Neither unity nor plurality should be extended
beyond the conditions under which it has been established.
A claim concerning one domain or determination cannot
automatically be generalized to every domain or
determination.

The chapter must therefore establish not merely that
unity and plurality coexist, but the conditions under
which particular forms of unity and plurality can
be attributed to particular ontological subjects.

\subsection{Methodological Requirements for the Inquiry}

The subsequent examination of unity and plurality
will proceed according to several methodological
requirements.

\begin{enumerate}
    \item \textbf{Specification of the subject:}
    Every claim of unity or plurality must identify
    the existent, identity, composition, aggregate,
    or other ontological subject under consideration.

    \item \textbf{Specification of the determination:}
    The analysis must identify whether unity or plurality
    concerns identity, composition, structure, relation,
    manifestation, or another determination.

    \item \textbf{Non-reduction:}
    Unity must not be reduced in advance to singularity,
    aggregation, structural coherence, or invariance.
    Plurality must not be reduced in advance to the
    plurality of components or identities.

    \item \textbf{Respect for ontological distinctions:}
    Distinctions established in the preceding chapters,
    especially those concerning identity, determination,
    composition, and invariance, must remain operative.

    \item \textbf{Separation of established claims and open questions:}
    A conceptual distinction does not automatically
    establish a complete ontological theory. Where
    the conditions of a claim remain unresolved, the
    question must remain open.

    \item \textbf{No premature universalization:}
    The analysis must not assume that every existent,
    every composition, or every domain exhibits one
    and the same form of unity or plurality.
\end{enumerate}

These requirements establish the conditions under
which the chapter can proceed without confusing
different ontological levels or treating a limited
claim as a universal principle.

\subsection{Scope of the Present Section}

The present section has formulated the ontological
problem of unity and plurality and identified the
principal distinctions required for its examination.

It has not yet established a complete definition
of unity, a universal account of plurality, or a
final criterion for determining the number of
identities associated with an existent.

Nor has it established that unity is always prior
to plurality, that plurality is always derived
from unity, or that one can be reduced to the other.

The following sections will examine these questions
in a determinate order.

Section~\ref{sec:unity-of-identity} investigates
the unity of identity.

Section~\ref{sec:unity-of-source-and-aggregate}
examines the distinction between unity of source
and unity of aggregate.

Section~\ref{sec:plurality-of-identities} addresses
the possibility and ontological status of plurality
of identities.

Section~\ref{sec:plurality-and-composition} examines
the relation between plurality and composition,
while Section~\ref{sec:identity-and-multiplicity}
investigates the broader relation between identity
and multiplicity.

Finally, the chapter distinguishes unity from
singularity, establishes the limits of claims
concerning unity and plurality, and synthesizes
the results of the inquiry.

The governing methodological commitment is that
unity and plurality must be understood through
their respective ontological subjects and
determinations. Neither can be defined adequately
by reducing it to the other, and neither can be
treated as a universal explanation of the
ontological constitution of existents.
